\section{The original fourth coupling coefficient as an arithmetic heat automorphism}
\label{sec:x4-arithmetic}

This section uses the complete finite-box Wilson calculation retained in
the complete source proof in Appendix~\ref{app:xi4-source}. Its parameter is
\(\xi=b/\kappa=1/(4g^4)\), whereas the completed Riemann function is
\(\xi(\rho)\) and the unchanged arithmetic heat function is \(Z_t(s)\).
These notations retain their original meanings. We construct the actual
operator-valued arithmetic observable, prove its fourth coupling
coefficient, and give the inverse of the resulting arithmetic map.

\subsection{Original source operators and the complete spectral functional}

Fix the original open box with vertices \(\{-L,\ldots,L\}^3\),
integer \(L\geq2\), its \(N=3(2L)(2L+1)^2\) positively oriented
physical edges, and its \(M=3(2L)^2(2L+1)\) contained elementary
faces. On the gauge-invariant subspace of product Haar probability
\(L^2(SU(2)^N)\), retain
\begin{gather}
T_a=-i\sigma_a/2,\qquad E_e=-\sum_{a=1}^3X_{e,a}^2,\qquad
H_0=\sum_eE_e,\qquad S=\sum_pW_p,\notag\\
H=\kappa H_0+b\sum_p(2-W_p)
 =\kappa(H_0-\xi S)+2bMI,\notag\\
\kappa=\frac{2g^2}{a},\quad b=\frac1{2g^2a},\quad
\xi=\frac1{4g^4},\quad a,g>0.
\label{xat:original-H}
\end{gather}
Here \(X_{e,a}\) is the original left derivative generated by
\(T_a\), and \(W_p\) is the original oriented fundamental Wilson
trace. The metric identity is \(-\Delta_e^c=4E_e\), so the electric
coefficient in that metric is \(\kappa/4\). Neither this factor nor
the scalar \(2bM\) is changed. Let \(\psi_\xi\) be the positive
unit ground vector and write
\begin{gather}
\mathcal E(\xi)=2bM+\kappa e(\xi),\qquad
A_\xi=\frac{H-\mathcal E(\xi)}{\kappa}
      =H_0-\xi S-e(\xi)I,\notag\\
\gamma_e(\xi)=\langle\psi_\xi,E_e\psi_\xi\rangle,\qquad
v_e(\xi)=(E_e-\gamma_e(\xi))\psi_\xi.
\label{xat:excitation}
\end{gather}
The inner product is conjugate-linear in its first entry. For real
\(\xi\) near zero, \(A_\xi\) is self-adjoint on the physical
\(H^2\) domain, nonnegative, and has the original smooth ground
vector in its kernel. The formula for \(A_\xi\) is the exact
physical excitation energy divided by its original \(\kappa\);
the scalar in \(H\) cancels its identical scalar in
\(\mathcal E\) in this displayed difference.

For distinct edges \(e,f\) sharing no face, let \(a_{ef}\) count
the ordered-by-containment pairs of adjacent faces containing
\(e\) and \(f\), respectively. The complete spectral coefficient
proved in the source calculation is, for bounded continuous complex
\(h:\R\to\C\) differentiable at \(3\),
\begin{equation}
\lim_{\xi\to0}\xi^{-4}
 \langle v_e(\xi),h(A_\xi)v_f(\xi)\rangle
 =a_{ef}\mathcal D(h),
\qquad
\mathcal D(h)=c_0h(3)-\frac{h'(3)}{336}
                  +c_1h(9/2)+c_2h(13/2),
\label{xat:source-functional}
\end{equation}
where every coefficient is retained:
\begin{equation}
c_0=-\frac{59}{275184},\qquad
c_1=\frac1{1296},\qquad c_2=\frac3{132496}.
\label{xat:coefficients}
\end{equation}
Equivalently,
\(\mathcal D=c_0\delta_3+\delta'_3/336+
c_1\delta_{9/2}+c_2\delta_{13/2}\), with
\(\delta'_3(h)=-h'(3)\). Thus the derivative term is part of the
original spectral coefficient. The source proof obtains it from the
actual first interacting spectral band and its full Gram isometry,
together with both exterior spin channels; it is not replaced here
by an atomic measure.

\subsection{Bounded arithmetic functional calculus at every complex heat time}

\begin{lemma}[The actual arithmetic test and its operator integral]
\label{xat:bounded-test}
For every complex \(t,s\), the function
\(h_{t,s}(q)=Z_t(s+q)\) is bounded and continuously differentiable
on the whole real \(q\)-axis. On real \(\xi\) near zero, its
bounded spectral calculus is exactly
\begin{equation}
Z_t(s+A_\xi)
 =4\int_\R k(v)e^{tv^2/4+isv}e^{ivA_\xi}\,dv.
\label{xat:operator-Fourier}
\end{equation}
The integral is the strong integral of bounded operators, defined on
each Hilbert-space vector; its operator norm is bounded by
\begin{equation}
4\int_\R |k(v)|
       e^{(\operatorname{Re}t)v^2/4+|\operatorname{Im}s||v|}\,dv.
\label{xat:operator-bound}
\end{equation}
It is jointly entire in \(t,s\), in the operator norm, and obeys
\(\partial_t Z_t(s+A_\xi)=-\tfrac14\partial_s^2Z_t(s+A_\xi)\).
\end{lemma}

\begin{proof}
Use the unchanged arithmetic identity \eqref{hc:Fourier}:
\[
h_{t,s}(q)=4\int_\R k(v)e^{tv^2/4+isv}e^{iqv}\,dv.
\]
For real \(q\), the last factor has modulus one. The previously
proved two-ended superexponential bound on \(k\) dominates the
absolute integral in \eqref{xat:operator-bound}. Its product with
\(|v|^j\) is also integrable for every nonnegative integer \(j\).
This proves boundedness of every real \(q\)-derivative, uniformly
in \(q\), by differentiating the integral. On compact complex
\((t,s)\) sets the same dominator works after every specified
\(t,s,q\) derivative.

For any Hilbert-space vector \(u\), the function
\(v\mapsto e^{ivA_\xi}u\) is strongly continuous and has constant
norm. Thus the right side of \eqref{xat:operator-Fourier} is a
Bochner integral on that vector and satisfies the stated norm bound.
For any two vectors, the corresponding spectral measure has finite
total variation, bounded by the product of their norms. Fubini
therefore interchanges this measure integral with the absolute
\(v\) integral and gives its matrix element as that of
\(h_{t,s}(A_\xi)\). This proves the operator identity, rather than
postulating an unbounded entire functional calculus.

Expand the scalar exponential about any fixed \((t,s)\). On every
smaller compact polydisc, its absolutely summed Taylor series is
dominated by the same integral with larger fixed quadratic and
linear weights. The integral norm inequality proves convergence
of the resulting operator Taylor series in norm. This establishes
joint entire dependence. Differentiation in \(t\) inserts
\(v^2/4\), whereas two \(s\) derivatives insert \(-v^2\), which
proves the heat identity with the original sign and factor.
\end{proof}

\begin{theorem}[The actual arithmetic fourth coefficient]
\label{xat:arithmetic-coefficient}
Define the entire arithmetic function
\begin{equation}
\mathcal K_t(s)=c_0Z_t(s+3)-\frac{Z_t'(s+3)}{336}
                 +c_1Z_t(s+9/2)+c_2Z_t(s+13/2).
\label{xat:K}
\end{equation}
For every complex \(t,s\), the actual original observable
\begin{equation}
G_\xi(t,s)=
 \langle v_e(\xi),Z_t(s+A_\xi)v_f(\xi)\rangle
\label{xat:G}
\end{equation}
satisfies
\begin{equation}
\lim_{\xi\to0}\xi^{-4}G_\xi(t,s)=a_{ef}\mathcal K_t(s).
\label{xat:limit}
\end{equation}
The limit retains the actual Riemann initial datum:
\begin{align}
\mathcal K_0(s)&=c_0\Xi(s+3)-\frac{\Xi'(s+3)}{336}
                    +c_1\Xi(s+9/2)+c_2\Xi(s+13/2),
\label{xat:initial}\\
\partial_t\mathcal K_t&=-\tfrac14\partial_s^2\mathcal K_t,
\qquad \partial_tG_\xi=-\tfrac14\partial_s^2G_\xi.
\label{xat:heat-identities}
\end{align}
\end{theorem}

\begin{proof}
Lemma~\ref{xat:bounded-test} verifies every test-function hypothesis
of \eqref{xat:source-functional} for the actual
\(h(q)=Z_t(s+q)\). Its values and derivative at the four specified
source locations give \eqref{xat:K} and \eqref{xat:limit}
literally. Every translated term and its retained derivative is
jointly entire. Translation by each of the original constants and
\(\partial_s\) commute with \(\partial_s^2\), so
\eqref{hc:heat} proves the first identity in
\eqref{xat:heat-identities}; the operator identity from
Lemma~\ref{xat:bounded-test} proves the second one. Finally
\(Z_0=\Xi\) gives \eqref{xat:initial} with no change of arithmetic
initial function.
\end{proof}

\subsection{Actual coupling analyticity and the sixth-order remainder}

\begin{proposition}[Uniform remainder on compact arithmetic parameter sets]
\label{xat:remainder}
For each original finite box and each compact
\(\mathcal Q\subset\C^2_{t,s}\), there are numbers
\(r>0\) and \(C_{\mathcal Q,L,e,f}<\infty\) such that, for real
\(|\xi|\leq r/2\),
\begin{equation}
\sup_{(t,s)\in\mathcal Q}
 |G_\xi(t,s)-a_{ef}\xi^4\mathcal K_t(s)|
 \leq C_{\mathcal Q,L,e,f}|\xi|^6.
\label{xat:uniform-remainder}
\end{equation}
The construction gives a jointly holomorphic extension in a complex
coupling disc and all complex \(t,s\). The radius and constant
retain their dependence on the original finite box.
\end{proposition}

\begin{proof}
We first specify the analytic continuation of the vectors and bras.
The full free ground eigenvalue is simple and the next full-space
eigenvalue is \(3/4\); the original perturbation has
\(\|S\|=2M\). The resolvent on \(|z|=3/8\) has norm at most
\(8/3\), so the source Neumann construction supplies the analytic
rank-one vacuum projection in a coupling disc contained in
\(|\xi|<3/(16M)\). Write \(\chi(\xi)=P(\xi)1\). Complex
conjugation \(C\) on functions satisfies
\(CP(\xi)C=P(\bar\xi)\), since both \(H_0\) and \(S\) have
real coefficients. Thus
\(C\chi(\xi)=\chi(\bar\xi)\).

The scalar
\(\mathfrak n(\xi)=\langle\chi(\bar\xi),\chi(\xi)\rangle\) is
holomorphic and equals one at zero. On a smaller symmetric disc,
it is nonzero and has the holomorphic square root with value one
at zero. Set \(\psi(\xi)=\chi(\xi)/\sqrt{\mathfrak n(\xi)}\).
For small real \(\xi\) this is the original positive unit vacuum:
it is the real normalized eigenvector with positive overlap with
the constant vector. For complex \(\xi\), this formula fixes its
analytic extension; no assertion that its Hilbert norm remains one
off the real axis is needed. Its eigenvalue \(e(\xi)\) is
holomorphic and real on the real axis.

Analyticity of \(\psi\) holds in every fixed Sobolev norm. Indeed
the exact equation
\[
\psi(\xi)=(I+H_0)^{-1}
             [(1+e(\xi)+\xi S)\psi(\xi)]
\]
increases the Sobolev index by two, since multiplication by the
original smooth \(S\) preserves that index and the elliptic inverse
increases it by two. Starting with the Hilbert-space analytic
projection proves the assertion inductively. Consequently
\[
\gamma_e(\xi)=
 \langle\psi(\bar\xi),E_e\psi(\xi)\rangle,
\qquad v_e(\xi)=(E_e-\gamma_e(\xi))\psi(\xi)
\]
are holomorphic in every required Sobolev space. Their restrictions
to real \(\xi\) are exactly \eqref{xat:excitation}. The
conjugate-linear-first functional
\(u\mapsto\langle v_e(\bar\xi),u\rangle\) is holomorphic in
\(\xi\); it is the precise continuation of the original bra.

On the same physical domain put
\(A_\xi=H_0-\xi S-e(\xi)I\), now for complex \(\xi\).
The bounded perturbation of the free unitary group gives
\(e^{ivA_\xi}\) for every real \(v\). Its Dyson series is
obtained by inserting \(-\xi S\) between the unchanged free
groups. The term with \(n\) insertions has norm at most
\((2M|\xi||v|)^n/n!\); the ordered integrals are strong
integrals on each vector. The resulting bounded-operator series
converges in norm, uniformly on compact coupling sets. The scalar
shift commutes with every term and contributes
\(e^{-iv e(\xi)}\). Hence
\begin{equation}
\|e^{ivA_\xi}\|
 \leq\exp\{ |v|(2M|\xi|+|\operatorname{Im}e(\xi)|)\}.
\label{xat:Dyson-bound}
\end{equation}
This proves norm-holomorphic dependence on \(\xi\) for fixed
real \(v\), with the entire original bounded potential retained.

Define the continuation of \eqref{xat:G} by
\begin{equation}
\widetilde G(\xi,t,s)=4\int_\R k(v)e^{tv^2/4+isv}
 \langle v_e(\bar\xi),e^{ivA_\xi}v_f(\xi)\rangle\,dv.
\label{xat:continued-G}
\end{equation}
Choose a closed coupling disc \(|\xi|\leq r\) strictly inside
the disc just constructed. Put
\[
B_r=\sup_{|\xi|\leq r}
 \|v_e(\bar\xi)\|\,\|v_f(\xi)\|,
\qquad
E_r=\sup_{|\xi|\leq r}|\operatorname{Im}e(\xi)|.
\]
Both numbers are finite by continuity on that closed disc. On a
compact arithmetic parameter set \(\mathcal Q\), set
\(T_{\mathcal Q}=\max_{\mathcal Q}\operatorname{Re}t\) and
\(S_{\mathcal Q}=\max_{\mathcal Q}|\operatorname{Im}s|\).
The absolute integral in \eqref{xat:continued-G} is uniformly
bounded by the explicit finite number
\begin{equation}
M_{r,\mathcal Q}=4B_r\int_\R |k(v)|
 e^{T_{\mathcal Q}v^2/4+
       (S_{\mathcal Q}+2Mr+E_r)|v|}\,dv.
\label{xat:Cauchy-bound}
\end{equation}
The superexponential arithmetic kernel absorbs all these retained
weights. On larger compact arithmetic sets it also absorbs the
weights introduced by each parameter derivative. The locally
uniformly convergent Dyson and scalar exponential series, or their
Cauchy integral formulas dominated by this same bound, therefore
prove joint holomorphy of \(\widetilde G\). For real coupling it
equals the original observable by
Lemma~\ref{xat:bounded-test}.

The source central-link unitary is
\[
(\mathcal ZF)(U)=
 F\bigl(((-1)^{\sum_{j<i}n_j}U_{(n,i)})_{(n,i)}\bigr).
\]
It preserves Haar measure, commutes with every original \(E_e\),
and changes the sign of every Wilson trace. Consequently
\(\mathcal Z(H_0-\xi S)\mathcal Z=H_0+\xi S\).
Uniqueness and positivity of the real vacuum imply
\(\psi_{-\xi}=\mathcal Z\psi_\xi\),
\(e(-\xi)=e(\xi)\),
\(v_e(-\xi)=\mathcal Z v_e(\xi)\), and
\(A_{-\xi}=\mathcal ZA_\xi\mathcal Z\).
Thus \(G_{-\xi}(t,s)=G_\xi(t,s)\) for real coupling. The identity
theorem extends this equality to the complex disc. No time or
arithmetic coordinate changes in this symmetry.

Write the coupling Taylor series as
\(\widetilde G=\sum_{n\geq0}g_n(t,s)\xi^n\).
The actual limit \eqref{xat:limit} on the real axis forces
\(g_0=g_1=g_2=g_3=0\) and
\(g_4=a_{ef}\mathcal K_t(s)\): otherwise division by \(\xi^4\)
could not have the proved finite limit. Evenness makes every odd
coefficient vanish. The Cauchy bound on \(|\xi|=r\) gives
\(\sup_{\mathcal Q}|g_n|\leq M_{r,\mathcal Q}r^{-n}\).
For \(|\xi|\leq r/2\), summing the remaining series yields
\[
\sup_{\mathcal Q}
 |\widetilde G-a_{ef}\xi^4\mathcal K|
 \leq 2M_{r,\mathcal Q}r^{-6}|\xi|^6.
\]
This proves \eqref{xat:uniform-remainder}, for example with the
specified constant \(2M_{r,\mathcal Q}r^{-6}\). All box dependence
in \(M,N,r,B_r,E_r\) remains present. The estimate does not
exchange a volume limit with a coupling derivative.
\end{proof}

\subsection{An invertible arithmetic map on the full retained Fourier domain}

Retain the exact spaces \(\E,\A\) and transform \(T\) already
defined above, including every seminorm
\[
p_{A,B,m,j}(\phi)=\int_\R|\phi^{(j)}(v)|(1+|v|)^m
                         e^{Av^2+B|v|}\,dv,
\qquad T\phi(s)=4\int_\R\phi(v)e^{isv}\,dv.
\]
Define the source multiplier and arithmetic operator by
\begin{align}
m(v)&=c_0e^{3iv}-\frac{iv}{336}e^{3iv}
               +c_1e^{9iv/2}+c_2e^{13iv/2},
\label{xat:multiplier}\\
(\mathcal T_{\mathcal D}f)(s)
 &=c_0f(s+3)-\frac{f'(s+3)}{336}
                  +c_1f(s+9/2)+c_2f(s+13/2).
\label{xat:T-D}
\end{align}

\begin{lemma}[A uniform exact lower bound for the original multiplier]
\label{xat:nonzero-multiplier}
For every real \(v\),
\begin{equation}
|m(v)|\geq\delta,\qquad \delta=\frac{575}{1752192}>0.
\label{xat:delta}
\end{equation}
In particular its real-axis reciprocal exists everywhere, with no
branch choice or removed spectral point.
\end{lemma}

\begin{proof}
Factoring the nonzero phase gives the exact identity
\[
e^{-3iv}m(v)=c_0-\frac{iv}{336}
                  +c_1e^{3iv/2}+c_2e^{7iv/2}.
\]
For \(|v|\leq1/2\), the elementary inequality
\(\cos x\geq1-x^2/2\) yields
\begin{align*}
\operatorname{Re}(e^{-3iv}m(v))
 &=c_0+c_1\cos(3v/2)+c_2\cos(7v/2)\\
 &\geq c_0+\frac{23}{32}c_1-\frac{17}{32}c_2
 =\frac{575}{1752192}.
\end{align*}
For completeness, this cosine bound follows directly by integration:
for \(x\geq0\),
\(1-\cos x=\int_0^x\sin u\,du\leq\int_0^xu\,du\), using
\(\sin u\leq u\); the latter follows from
\(u-\sin u=\int_0^u(1-\cos w)\,dw\geq0\). Evenness gives the
negative case. Since the real part is positive, its displayed
lower bound also bounds the modulus.

For \(|v|\geq1/2\), the triangle inequality with the original
linear term gives
\begin{align*}
|m(v)|&\geq\frac{|v|}{336}-(|c_0|+c_1+c_2)\\
 &\geq\frac1{672}-\frac{10825}{10732176}
 =\frac{10291}{21464352}
 >\frac{575}{1752192}.
\end{align*}
All coefficients in these estimates are the original rational
coefficients. The two intervals cover the whole real axis and
prove \eqref{xat:delta}.
\end{proof}

\begin{theorem}[Exact arithmetic automorphism and heat conjugacy]
\label{xat:automorphism}
Multiplication by \(m\) is a continuous automorphism of \(\E\).
Its inverse is literal pointwise multiplication by \(1/m\) on
the unchanged real Fourier axis. The operator
\(\mathcal T_{\mathcal D}\) is a continuous automorphism of
\(\A\), with exact inverse
\begin{equation}
\mathcal T_{\mathcal D}^{-1}f
 =T\left(\frac{T^{-1}f}{m}\right).
\label{xat:inverse}
\end{equation}
For every complex \(t\),
\begin{gather}
\mathcal T_{\mathcal D}=TM_mT^{-1},\qquad
\mathcal T_{\mathcal D}U_t=U_t\mathcal T_{\mathcal D},\qquad
\mathcal T_{\mathcal D}^{-1}U_t
        =U_t\mathcal T_{\mathcal D}^{-1},
\label{xat:commutation}\\
\mathcal K_t=\mathcal T_{\mathcal D}Z_t,
\qquad
Z_t=\mathcal T_{\mathcal D}^{-1}\mathcal K_t.
\label{xat:actual-inverse}
\end{gather}
Thus the actual arithmetic function is exactly recoverable from
the propagated source coefficient, on this fully specified domain.
\end{theorem}

\begin{proof}
Every derivative of \(m\) has at most linear growth on the real
axis. More explicitly, for \(j\geq0\) it satisfies
\( |m^{(j)}(v)|\leq C_j(1+|v|)\), where one may take
\[
C_j=|c_0|3^j+c_1(9/2)^j+c_2(13/2)^j
       +\frac{3^j+\mathbf1_{j\geq1}j3^{j-1}}{336}.
\]
The product rule gives the seminorm estimate
\[
p_{A,B,m,j}(m\phi)
 \leq\sum_{\ell=0}^j\binom j\ell C_\ell
                     p_{A,B,m+1,j-\ell}(\phi).
\]
This proves continuity of \(M_m\) on the original \(\E\).

Let \(r(v)=1/m(v)\), which is smooth by
Lemma~\ref{xat:nonzero-multiplier}. The exact product identity
\(mr=1\) and its derivatives give
\[
r^{(j)}=-\frac1m\sum_{\ell=1}^j
                 \binom j\ell m^{(\ell)}r^{(j-\ell)}
\quad(j\geq1),\qquad |r|\leq\delta^{-1}.
\]
Set \(Q_0=\delta^{-1}\) and recursively define the finite
positive constants
\[
Q_j=\delta^{-1}\sum_{\ell=1}^j
                      \binom j\ell C_\ell Q_{j-\ell}.
\]
Induction proves
\(|r^{(j)}(v)|\leq Q_j(1+|v|)^j\): the term with index
\(\ell\geq1\) has exponent at most
\(1+j-\ell\leq j\), and the denominator is bounded by
\(\delta^{-1}\). Therefore
\[
p_{A,B,m,j}(r\phi)
 \leq\sum_{\ell=0}^j\binom j\ell Q_\ell
                         p_{A,B,m+\ell,j-\ell}(\phi).
\]
This proves that reciprocal multiplication is continuous on
exactly the same \(\E\). The two pointwise compositions are
the identity, proving the Fourier automorphism and its inverse.

Applying the retained transform to each term of \(m\phi\)
gives the three translations and the derivative in
\eqref{xat:T-D}; the factor \(iv\) is exactly the Fourier
multiplier of \(\partial_s\). The injectivity of \(T\) already
proved above makes its inverse on \(\A\) unambiguous. This
proves the first identity in \eqref{xat:commutation} and
\eqref{xat:inverse}. Multiplication by \(m\) and by \(1/m\)
each commutes pointwise with multiplication by
\(e^{tv^2/4}\). All three maps preserve \(\E\), so applying
\(T\) proves the two full-domain heat commutation identities
for every complex \(t\). Finally \eqref{xat:K} is exactly
\(\mathcal T_{\mathcal D}Z_t\), and applying the proved inverse
gives \eqref{xat:actual-inverse}.
\end{proof}

The source factor acts by the typed scalar map
\(M_{a_{ef}}:\A\to\A,\ f\mapsto a_{ef}f\).
For \(a_{ef}>0\) its inverse is \(g\mapsto g/a_{ef}\), by direct
composition; for \(a_{ef}=0\) its kernel is all of \(\A\) and its
image is \(\{0\}\). This retains the complete pair observation map
\(M_{a_{ef}}\mathcal T_{\mathcal D}\), including its zero case.
When \(a_{ef}>0\), its division also recovers the unscaled
\(\mathcal K\) from the coefficient in
\eqref{xat:limit}. The original boxes contain such pairs. In
particular, for \(e=((0,0,0),1)\) and
\(f=((0,2,0),1)\), there is exactly one common edge-adjacency
neighbor, \(((0,1,0),1)\). To verify uniqueness directly, a
direction-1 neighbor of \(e\) is its parallel translate by one
step in direction 2 or 3; the corresponding lists for \(e\)
and \(f\) meet only at the displayed midpoint edge. A
direction-2 neighbor has its direction-2 base coordinate in
\(\{-1,0\}\) for \(e\), and in \(\{1,2\}\) for \(f\),
so these lists are disjoint. A direction-3 neighbor has
direction-2 coordinate zero for \(e\) and two for \(f\), also
disjoint. All faces in the unique midpoint channel are contained
in every original box with \(L\geq2\). The exact source
channel bijection gives \(a_{ef}=1\), including these open
boundaries.

Equations \eqref{xat:continued-G} and \eqref{xat:actual-inverse}
provide the complete source-to-arithmetic heat correspondence.
They assert an invertible map of the specified arithmetic function space.
The coordinate-level map on every pointwise zero jet and analytic unit
is proved next.
The exact morphism on zero data uses the transported jet maps below.
Fix $z\in\C$ and an integer $n\ge1$. Define continuous linear maps
\begin{align*}
\mathcal J_{z,n}:\A&\longrightarrow\C^n,
&\mathcal J_{z,n}f&=(f(z),f'(z),\ldots,f^{(n-1)}(z)),\\
\mathcal L_{z,n}:\A&\longrightarrow\C^n,
&\mathcal L_{z,n}g&=\mathcal J_{z,n}\mathcal T_{\mathcal D}^{-1}g.
\end{align*}
For $g=T\phi$ their individual coordinate functionals are exactly
\begin{equation}
(\mathcal L_{z,n}g)_j
 =4\int_\R (iv)^j\frac{\phi(v)}{m(v)}e^{izv}\,dv,
\qquad j=0,\ldots,n-1.
\label{x4:transported-jets}
\end{equation}
This follows by differentiating the proved inverse integral; it retains
the factor $4$, the original coordinate $z$ and every derivative factor
$(iv)^j$. Absolute values are bounded by
\[
4\delta^{-1}\int_\R|\phi(v)||v|^j e^{|\operatorname{Im}z||v|}\,dv,
\]
where $\delta=575/1752192$. One defining seminorm with an integer
linear weight at least $|\operatorname{Im}z|$ dominates this integral.
The same estimate on compact sets justifies every displayed derivative
and proves continuity. It also proves continuity of $\mathcal J_{z,n}$
by omitting the factor $1/m$.

Both maps are onto. To prove this without assuming a jet-extension
theorem on $\A$, set $u_z(s)=\Xi(s-z)$. This lies in $\A$, since its
kernel is $k(v)e^{-izv}$ and all differentiated weighted seminorms
are bounded by finitely many seminorms of $k$ with a larger linear
weight. Each $(s-z)^ju_z(s)$ also lies in $\A$ by the already proved
invariance under multiplication by $s$. Its first $n$ derivatives at
$z$ form a triangular matrix with diagonal entries
$j!\Xi(0)\ne0$, $j=0,\ldots,n-1$. The nonzero value was proved by
the positive kernel integral. Back substitution therefore constructs
every specified jet exactly. This proves surjectivity of
$\mathcal J_{z,n}$, and that of $\mathcal L_{z,n}$ follows by
composition with the automorphism.
The finite back-substitution formula is also a continuous linear right
inverse from $\C^n$ into the displayed finite-dimensional span.

Put $\mathcal M_{z,n}=\ker\mathcal J_{z,n}$. The exact zero-jet
correspondence is
\[
\mathcal T_{\mathcal D}\mathcal M_{z,n}
 =\ker\mathcal L_{z,n},\qquad
\A/\mathcal M_{z,n}\longrightarrow
\A/\ker\mathcal L_{z,n},\quad
[f]\longmapsto[\mathcal T_{\mathcal D}f].
\]
The first equality follows in both directions by applying the inverse
automorphism. It makes the quotient map well-defined and gives the
inverse $[g]\mapsto[\mathcal T_{\mathcal D}^{-1}g]$; both are
continuous by their quotient topologies. Both quotients have the
displayed coordinates in $\C^n$, and the map is the identity in those
coordinates because
$\mathcal L_{z,n}\mathcal T_{\mathcal D}=\mathcal J_{z,n}$.

For the actual $\mathcal K_t=\mathcal T_{\mathcal D}Z_t$ this proves
every jet identity
$\mathcal L_{z,n}\mathcal K_t=\mathcal J_{z,n}Z_t$.
Thus a zero of order exactly $n$ of $Z_t$ is expressed on the transformed
function by vanishing of coordinates $0,\ldots,n-1$ of
$\mathcal L_{z,n+1}\mathcal K_t$ and nonvanishing of coordinate $n$.
At every such zero the complete analytic unit is recovered as the
convergent series
\[
\frac{Z_t(s)}{(s-z)^n}
 =\sum_{j=0}^\infty
 \frac{(\mathcal L_{z,n+j+1}\mathcal K_t)_{n+j}}{(n+j)!}(s-z)^j.
\]
It is the Taylor expansion of the recovered entire function, so the
identity and convergence follow from its holomorphy. No derivative,
unit, coordinate or multiplicity is discarded by this morphism.

